\makeatletter
\providecommand{\equationalias}[1]{\protected@write\@auxout{}{\string\newlabel{#1}{{\theequation}{\thepage}{}{equation.\theHequation}{}}}}
\makeatother

\section{The result}\label{sec:introduction}\label{sec:result}

How can a lattice model give a quantum field theory in continuous space-time? We answer this question for four-dimensional Yang--Mills theory with a compact gauge group whose Lie algebra is simple and nonabelian, starting from Wilson's lattice gauge theory~\cite{Wilson1974}. On a finite periodic lattice the variables are group elements attached to the links, and their law is an ordinary probability measure. Let $G\subset U(n_G)$ be a closed matrix group, $s\ge1$ an integer and $\Lambda_s=(\Z/s\Z)^4$. A configuration assigns $U(y,a)\in G$ to the positive link from the site~$y$ in the coordinate direction~$a$, and $U(y,a)^{-1}$ to the reversed link. For $\beta\in\R$,
\begin{equation}\label{eq:intro-wilson}
\begin{aligned}
U_{y;ab}&=U(y,a)U(y+e_a,b)U(y+e_b,a)^{-1}U(y,b)^{-1},\\
E_s^{\rm W}(U)&=\sum_{y\in\Lambda_s}\sum_{a<b}\bigl(1-n_G^{-1}\operatorname{Re}\tr U_{y;ab}\bigr),\qquad
d\mu_{s,\beta}(U)=Z_{s,\beta}^{-1}e^{-\beta E_s^{\rm W}(U)}\prod_{y,a}dh(U(y,a))\,,
\end{aligned}
\end{equation}
where $dh$ is probability Haar measure and $Z_{s,\beta}$ normalizes the measure. The plaquette holonomy $U_{y;ab}$ is the transport around an elementary square, and $\beta$ penalizes its distance from the identity. The problem is to choose the inverse coupling and the local observables as functions of the lattice spacing so that their correlations converge as the spacing tends to zero and the volume grows, and to identify the limit.

\paragraph{Main result.}
Let the real Lie algebra $\g$ of~$G$ be simple and nonabelian; $G$ need not be connected. Fix a block penalty $\blockpenalty>0$: the coefficient of the soft constraint that ties each coarse link to a predictor computed from the fine transports in the block-spin kernel (Definition~\ref{def:D04}). We construct one family of Wilson measures on tori of physical side $L_p\uparrow\infty$ with spacings $a_N=D_N^{-1}\downarrow0$, together with continuous gauge-invariant local insertions on fixed finite stencils, all chosen before any sequence of volumes is prescribed. Its bare inverse couplings satisfy the one-loop running law
\begin{equation}\label{eq:intro-running}
\frac{\beta_{N,p}}{\log(1/a_N)}\longrightarrow\slope=\frac{11\casimir}{24\pi^2}
\qquad(N\to\infty),\quad\text{uniformly in }p\,,
\end{equation}\equationalias{eq:I12-bare-slope}%
where $\casimir$ is the adjoint Casimir in the trace normalization of Definition~\ref{def:D01} (for $SU(n)$, $\slope=11n^2/(12\pi^2)$; Remark~\ref{rem:SUn}). The family has the following properties; Theorem~\ref{thm:B} in Section~\ref{sec:main-theorem} states them precisely.
\begin{enumerate}[label=\textup{(\roman*)},nosep]
\item At each fixed volume, correlations of the insertions, smeared against smooth tests supported away from coincident points and inside one period cell, converge along the full sequence of spacings.
\item Every prescribed strictly increasing sequence $\mathsf p$ of volume indices has a strictly increasing refinement $\mathsf q$ such that, with $P=\mathsf p\circ\mathsf q$ and the lattice correlations $\latsmear_{N,p}$ of Definition~\ref{def:D03},
\begin{equation}\label{eq:intro-ordered-limit}
\lim_{k\to\infty}\Bigl(\lim_{N\to\infty}\latsmear_{N,P(k)}(v;f)\Bigr)=\int f(x)S_n(v;x)\,dx
\end{equation}
for one table $S$ of Euclidean correlation functions, symmetric under permutations, translations and signed permutations of the axes, reflection positive, smooth away from collisions and polynomially bounded near them (properties E1--E7 of Section~\ref{sec:euclidean-properties}). One refinement serves all insertions and all tests.
\item $S=\int_{\phases}S^\xi\,d\nu(\xi)$ is an average of phases over a compact metric space with a fully supported probability measure, and each $S^\xi$ is the Euclidean continuation of a faithful Wightman theory with unique vacuum.
\item The fields carry a conserved symmetric stress tensor, normalized by its action on the other fields. In a selected phase, the two-point function $\tracetwopoint$ of the centered trace of this tensor satisfies
\begin{equation}\label{eq:intro-trace-law}
\lim_{r\downarrow0}r^8\log^2(1/r)\,\tracetwopoint(r)=\frac{3\dim_\R\g}{\pi^4}>0\,,
\end{equation}
and on no interval $(0,\delta)$ does $\tracetwopoint$ agree with a finite Wick profile, a finite sum of finite products of free mass-shell Laplace kernels (Definition~\ref{def:D17}).
\item Nearby fields have an operator product expansion whose coefficients depend neither on the other fields nor on the phase, with remainders bounded uniformly in the phase. The correlations of each phase extend across coincident points by an explicit subtraction of Taylor jets.
\end{enumerate}
Theorem~\ref{thm:B} holds for every $\blockpenalty>0$. Appendix~\ref{sec:correspondence} states the existence theorem in the exact form of the formal statement (Theorem~\ref{thm:A}), which asks only for some positive logarithmic slope, and deduces it from Theorem~\ref{thm:B}.

\paragraph{What is not claimed.}
The Clay problem~\cite[\S4]{JaffeWitten2006} asks for existence and a positive mass gap. The construction gives Wightman theories with unique vacuum, one in each phase; no mass gap is proved or used, and uniqueness of the vacuum gives no lower bound on the rest of the energy spectrum. The infinite-volume limit, the phase space and the selected phase are not shown to be unique, and independence from~$\blockpenalty$ is not proved. Section~\ref{sec:main-theorem} lists the limitations next to Theorem~\ref{thm:B}.

\subsection{Source of the result and of this text}\label{sec:intro-source}

% BEGIN-VERIFICATION-PARAGRAPH
The mathematical result is due to OpenAI's Lean construction at commit~\texttt{3456e5e} of the repository \href{https://github.com/yang-mills-group/OpenAI-YM-20261028}{\texttt{yang-mills-group/OpenAI-YM-20261028}}~\cite{OpenAI2026YM}. Theorem~\ref{thm:B} records that construction; Appendix~\ref{sec:correspondence} states its top-level existence theorem and deduces it. Three kinds of checking are involved, and they should be kept apart.
\begin{itemize}[nosep]
\item The Lean kernel checks the formal proof. Its \href{https://github.com/yang-mills-group/OpenAI-YM-20261028/blob/3456e5e75cf8d4cd558955a5c7638fba2d68f464/verification/joint-check.log}{kernel-check record} reports no unproved placeholders in the proved theorem and only propositional extensionality, classical choice and quotient soundness as axioms.
\item AI agents compared every numbered statement of this text with the Lean declarations used by the main theorem: hypotheses, definitions, quantifier order, normalizations and conclusions, with usage determined from the compiled dependencies. The correspondence table is published separately; the source repository supplies a \href{https://github.com/yang-mills-group/OpenAI-YM-20261028/blob/3456e5e75cf8d4cd558955a5c7638fba2d68f464/verification/declaration-usage.tsv}{declaration-usage table}.
\item The prose arguments were written and reviewed by AI agents. They are not machine-checked and have not been refereed by human mathematicians.
\end{itemize}
A correct formal declaration does not by itself show that an English statement has kept its hypotheses and quantifiers, or that a prose proof proves that statement. This text was written entirely by AI agents under human direction, with separate writing and checking of statements and arguments, by the method described in \href{https://github.com/scottnarmstrong/yang-mills-informalization/blob/main/manuscript/PROCESS.md}{\emph{How the manuscript is written}}. The work is carried out in collaboration with Martin Hairer's \href{https://github.com/hairer/humanise-yang-mills}{\emph{humanise-yang-mills}} project.
% END-VERIFICATION-PARAGRAPH

\subsection{Earlier work and the method}\label{sec:intro-earlier}

Table~\ref{tab:earlier-work} lists, for the main mechanisms of the proof, the closest earlier rigorous result and its use here. The comparison concerns how these operations are combined; it assigns none of them a new origin, and we make no priority claim for renormalization-group principles or physical predictions.

\begingroup
\small
\setlength{\tabcolsep}{4pt}
\renewcommand{\arraystretch}{1.12}
\begin{longtable}{@{}p{0.25\textwidth}p{0.42\textwidth}p{0.28\textwidth}@{}}
\caption{Mechanisms, earlier results and their use here.}\label{tab:earlier-work}\\
\toprule
Mechanism & Earlier result & Use here\\
\midrule
\endfirsthead
\toprule
Mechanism & Earlier result & Use here\\
\midrule
\endhead
\bottomrule
\endfoot
Nonlinear block averages, gauge fixing, stationary backgrounds, marginal curvature-square extraction & Ba\l{}aban: analytic stationary backgrounds~\cite[Theorem~1, Proposition~6]{Balaban1985Variational}; symmetric block kernels and the marginal analysis~\cite[pp.~252--254, \S\S4--5]{Balaban1987RG}; large fields for small effective couplings~\cite[Theorem~1, p.~355]{Balaban1989Large}, physical observables left open there (p.~356) & Running couplings and local composite fields controlled together on one tuned family\\[2pt]
Four-dimensional ultraviolet construction & Magnen--Rivasseau--S\'en\'eor~\cite{MagnenRivasseauSeneor1993}: pure $SU(2)$, trivial sector, regularized axial gauge, fixed infrared cutoff; detailed convergence deferred and infrared cutoff retained (p.~327), composite observables set aside (p.~329) & Infinite-volume gauge-invariant local fields, phases, stress tensor and products\\[2pt]
Disjoint polymers covering the large-field region; selective exponentiation & Gaw\k{e}dzki--Kupiainen, critical $\varphi^4_4$~\cite[p.~212, (4.8), pp.~218--221]{GawedzkiKupiainen1985} & Forced activities covering all large-field sites\\[2pt]
Polymer extraction, reblocking, analytic norms, Gaussian integration & Brydges--Yau~\cite[\S\S1--2, pp.~378--384]{BrydgesYau1990} & Activities and their weighted norms\\[2pt]
Clipped kinetic term, local quadratic terms, regular errors, covering sum, exact coupling reset & OpenAI's sigma-model papers: the exact map of~\cite[\S5, pp.~38--55]{OpenAI2026O4} and its adaptation~\cite[\S3.3, pp.~23--28]{OpenAI2026O3} & The coordinates of the effective densities\\[2pt]
Reflection positivity, transfer matrix, chessboard estimates & Osterwalder--Seiler~\cite{OsterwalderSeiler1978}, \cite[\S2a, Theorem~2.1]{Seiler1982}; Fr\"ohlich--Israel--Lieb--Simon~\cite[\S4, Theorems~4.1--4.2]{FrohlichIsraelLiebSimon1978}, lattice version for local functions~\cite[Theorem~2.2]{Seiler1982} & Volume-uniform bounds; reflection Hilbert spaces\\[2pt]
Analytic continuation of Euclidean correlations & Osterwalder--Schrader~\cite{OsterwalderSchrader1973}, with the growth and boundary-value issues addressed in~\cite{OsterwalderSchrader1975} & Reconstruction in every phase, from slab and distribution estimates; reflection positivity alone does not give it\\[2pt]
One-loop running & Gross--Wilczek, Politzer~\cite{GrossWilczek1973,Politzer1973}; background-field coefficient~\cite[\S3.3]{Abbott1982} & The slope~$\slope$, on the exact block map\\[2pt]
Operator products, trace anomaly, Taylor subtraction & Wilson~\cite[\S II]{Wilson1969}; all-order perturbative OPE in Euclidean Yang--Mills theory with gauge-invariant closure and associativity~\cite{FrobHolland2016}; Collins--Duncan--Joglekar~\cite[(3.19)--(3.21)]{CollinsDuncanJoglekar1977}; Brunetti--Fredenhagen~\cite[\S5.2]{BrunettiFredenhagen2000} & A nonperturbative OPE at each finite truncation, uniform in the phase; the trace law; contact extensions\\
\end{longtable}
\endgroup

The renormalization group is organized as in the method of Bauerschmidt, Brydges and Slade~\cite{BauerschmidtBrydgesSlade2019}: an explicitly defined map, a finite-dimensional flow of couplings, norms in which the remaining coordinate contracts, and a global trajectory with a tuned initial parameter. What requires proof here is the compatibility of the one-loop slope with the exact gauge block map, the coupling reset and the tuning to a common endpoint, followed by continuum identifications on one family: control of composite fields, restoration of rotations through stress-tensor fluxes, the trace law through all radii and product bounds uniform over phases.

\subsection{Reading guide}\label{sec:reading-guide}

Section~\ref{sec:definitions} defines the objects needed to state Theorem~\ref{thm:B}, which is Section~\ref{sec:main-theorem}. Section~\ref{sec:strategy} explains the route of the proof and the mechanisms that make it work, with a dependency map (Figure~\ref{fig:dependency}). Sections~\ref{sec:renormalization} and~\ref{sec:continuum-results} state the principal results, each with the idea of its proof, and Section~\ref{sec:deduction} deduces Theorem~\ref{thm:B}. Part~II gives the proofs. Table~\ref{tab:principal-results} lists the principal results; their tags appear in statement titles. Appendix~\ref{sec:constants} records the order in which constants are chosen. A first reading can skip the coordinates of the effective densities (Section~\ref{sec:density}), the marked decomposition (Section~\ref{sec:marked}) and Appendix~\ref{sec:constants}. The index of notation follows the appendices.

\begingroup
\small
\setlength{\tabcolsep}{4pt}
\begin{longtable}{@{}p{0.06\textwidth}p{0.55\textwidth}p{0.12\textwidth}p{0.21\textwidth}@{}}
\caption{Principal results and their proofs.}\label{tab:principal-results}\\
\toprule
Tag & Statement & Part~I & Proof in Part~II\\
\midrule
\endfirsthead
\toprule
Tag & Statement & Part~I & Proof in Part~II\\
\midrule
\endhead
\bottomrule
\endfoot
RG & Exact integration of one scale and coupling drift & Thm~\ref{prop:RG} & Sections~\ref{sec:partII-B}--\ref{sec:partII-D}\\
MT & Common endpoint coupling and running slope & Thm~\ref{prop:MT} & Section~\ref{sec:partII-E.1}\\
UB & Volume-uniform bounds on lattice differences & Prop~\ref{prop:UB} & Section~\ref{sec:partII-E.2}\\
UV & Full fixed-volume ultraviolet limit & Thm~\ref{prop:UV} & Section~\ref{sec:partII-E.3}\\
IV & Volume refinement and Euclidean properties & Thm~\ref{prop:IV} & Section~\ref{sec:partII-F}\\
FC & Inclusions, derivatives and Maxwell dimensions & Prop~\ref{prop:FC} & Lemma~\ref{lem:F2}\\
MC & Maxwell comparison at short distances & Prop~\ref{prop:MC} & Lemma~\ref{lem:G2.1}\\
ST & Conserved stress tensor and its shell fluxes & Thm~\ref{prop:ST} & Lemmas~\ref{lem:G3.1}--\ref{lem:G3.2}\\
SC & Differentiated comparison with stress insertions & Prop~\ref{prop:SC} & Lemma~\ref{lem:G3.3}\\
PH & Phase decomposition on the same fields & Thm~\ref{prop:PH} & Section~\ref{sec:partII-H.1}\\
RC & Wightman reconstruction in every phase & Thm~\ref{prop:RC} & Section~\ref{sec:partII-H.2}\\
TR & All-radii law for the centered stress trace & Thm~\ref{prop:TR} & Section~\ref{sec:partII-I.1}\\
NF & Exclusion of finite Wick profiles & Prop~\ref{prop:NF} & Lemma~\ref{lem:I2.1}\\
OP & Operator product expansion with error bounds & Thm~\ref{prop:OP} & Section~\ref{sec:partII-J.1}\\
CT & Extensions through collision diagonals & Thm~\ref{prop:CT} & Section~\ref{sec:partII-J.2}\\
\end{longtable}
\endgroup

\subsection{Notation and conventions}\label{sec:conventions}

Throughout, $\N=\{0,1,\ldots\}$, coordinate indices run from $1$ to $4$, $e_a$ is the $a$th coordinate vector, and time is the fourth coordinate. For the relativistic statements write
\begin{equation*}
\minksquare(x)=(x^4)^2-|\vec x|^2,\qquad
\forwardcone=\{\mathbf p:|\vec{\mathbf p}|\le\mathbf p^4\}\,.
\end{equation*}
Inner products are antilinear in the first variable, and spacetime integrals use Lebesgue measure. Correlations of arity zero are included; empty sums are zero and empty products are one. The group $\hypercubic$ consists of the signed permutation matrices, and $\mathsf s_a^{\rm refl}$ denotes reflection in the $a$th coordinate.
